% 원본 교안 앞에 추가한 사전 실습. Nord VS Code의 실제 배포 파일·화면.
% 캡처와 실행 범위: assets/vscode-01/README.md
\newcommand{\vscpicture}[2][]{%
  \par\vspace{0.12cm}\centering
  \includegraphics[width=\textwidth,height=4.0cm,keepaspectratio,#1]{assets/vscode-01/#2.png}\par}
\newcommand{\vscnote}[1]{\par\vspace{0.16cm}{\small #1}\par}
\section{01: VS Code 사전 시뮬레이션}
\begingroup
% 새 안내의 설명과 캡처가 고정 로고·목차 영역에 닿지 않게 여백을 확보한다.
\renewcommand{\small}{\fontsize{9}{11}\selectfont}

\begin{frame}[label=legacy-01]{01 · 논리 게이트 실습}
{\Large\color{UOSBlue}\bfseries VS Code에서 먼저 확인한다}\par\vspace{0.45cm}
\begin{enumerate}
\setlength{\itemsep}{0.25cm}
\item 새 창에서 실습 workspace를 연다.
\item AND·OR·XOR 소스와 테스트벤치를 읽는다.
\item 네 입력 조합을 시뮬레이션하고 파형을 해석한다.
\item 실험 전 레포트를 작성한 뒤 Vivado GUI를 연다.
\end{enumerate}
\vspace{0.35cm}
{\small 앞부분: 새로 작성한 VS Code 사전 실습\par
뒷부분: 레거시 교안의 Vivado 절차와 원본 화면}\par\vspace{0.3cm}
{\small 작성일 \materialdate\quad\presenter}
\end{frame}

\begin{frame}[label=legacy-01-contents]{01 · 실습 목차}
\begin{enumerate}
\setlength{\itemsep}{0.23cm}
\item \navlink{vsc01-prepare}{준비: 도구·새 창·workspace}
\item \navlink{vsc01-extensions}{확장 확인과 코드 열기}
\item \navlink{vsc01-truth}{예상값·RTL·테스트벤치 읽기}
\item \navlink{vsc01-run}{시뮬레이션 작업 실행과 PASS 확인}
\item \navlink{vsc01-wave}{VaporView: 파형 열기·신호 추가·값 읽기}
\item \navlink{vsc01-report}{실험 전 레포트와 문제 해결}
\item \navlink{legacy-01-vivado}{그다음: Vivado 프로젝트·시뮬레이션·FPGA}
\end{enumerate}
\vscnote{좌하단 \textbf{목차} 버튼은 언제든 이 페이지로 돌아온다.}
\end{frame}

\begin{frame}[label=vsc01-prepare]{시작 전에 준비할 것}
\small
\begin{itemize}
\setlength{\itemsep}{0.20cm}
\item Windows PC에 VS Code와 Vivado 시뮬레이션 도구가 설치되어 있어야 한다. 저장소도 PC에 내려받는다.
\item VS Code는 편집기다. 이번 사전 실습의 실제 시뮬레이터는 \textbf{XSim}, 파형 뷰어는 \textbf{VaporView}다.
\item \texttt{xvlog}, \texttt{xelab}, \texttt{xsim}의 실행 경로를 준비한다. \textbf{Vivado GUI나 .xpr를 먼저 열 필요는 없다.}
\item 설치가 안 되어 있으면 같은 배포 폴더의\par
\texttt{02.vscode\_verilog\_환경설정\_매뉴얼.pdf}를 먼저 따른다.
\end{itemize}
\vspace{0.35cm}
이번 캡처 검증 환경: XSim 2026.1, slang 확장 0.2.17,\par
VaporView 1.5.4, 사용자의 Nord 테마.\par\vspace{0.2cm}
뒤의 레거시 화면과 도구 버전이 다르며, 레거시 Vivado·실제 보드의 새 실행 증빙을 뜻하지 않는다.
\end{frame}

\begin{frame}{1. VS Code 새 창부터 열기}
\small \textbf{File → New Window}를 누른다. 단축키는 \textbf{Ctrl+Shift+N}이다.
\vscpicture[trim=22.5bp 480bp 637.5bp 0bp,clip]{02-file-menu}
\vscnote{기존 창은 그대로 두고 실습용 창을 하나 더 연다.\par
한글 UI에서는 파일 → 새 창이다. Welcome 탭은 ×로 닫아도 된다.}
\end{frame}

\begin{frame}{2. workspace를 여는 메뉴 선택}
\small 새 창에서 \textbf{File → Open Workspace from File...}을 누른다.
\vscpicture[trim=22.5bp 352.5bp 637.5bp 0bp,clip]{02-file-menu}
\vscnote{workspace는 소스 폴더·편집 설정·시뮬레이션 작업을 함께 여는 파일이다.\par
\textbf{Open File...}로 Verilog 파일 하나만 여는 것과 구분한다.}
\end{frame}

\begin{frame}{3. 실습 01 workspace 선택}
\small 저장소에서 다음 폴더를 찾는다.\par
{\scriptsize\texttt{example/fpga\_projects\_hdl/LAB1/legacy/01\_logic\_gates}}
\vscpicture{03-select-workspace}
\vscnote{\texttt{legacy\_01\_logic\_gates.code-workspace}를 선택하고\par
오른쪽 아래 \textbf{Open}을 누른다. PC마다 저장소의 상위 경로는 다르다.}
\end{frame}

\begin{frame}{4. 올바른 창인지 확인}
\small 왼쪽 \textbf{Explorer}에서 \textbf{01\_LOGIC\_GATES}와 \textbf{COMMON}을 찾는다.
\vscpicture[trim=0bp 217.5bp 562.5bp 0bp,clip]{04-workspace}
\vscnote{\textbf{01\_LOGIC\_GATES}: 실행 작업·결과 폴더.\quad\textbf{COMMON}: RTL·테스트벤치.\par
목록이 안 보이면 왼쪽 맨 위 파일 모양 아이콘을 누른다.}
\end{frame}

\begin{frame}[label=vsc01-extensions]{5. 확장 두 개 확인}
\small
\begin{enumerate}
\setlength{\itemsep}{0.22cm}
\item 왼쪽 \textbf{Extensions} 아이콘(네모 모양)을 누른다.\par
단축키는 \textbf{Ctrl+Shift+X}다.
\item 검색창에 \texttt{@installed slang}을 입력한다.\par
\textbf{slang-server: Verilog/SystemVerilog LSP}를 선택한다.
\item 같은 방법으로 \texttt{@installed vaporview}도 찾는다.
\end{enumerate}
\vspace{0.25cm}
설치되어 있지 않으면 \texttt{@installed}를 빼고 검색한다.\par
배포자가 각각 \textbf{Hudson River Trading}, \textbf{Lloyd Ramseyer}인지 확인한다.
\vscnote{정확한 ID: \texttt{hudson-river-trading.vscode-slang}\par
\texttt{lramseyer.vaporview}. 설치 단계는 환경설정 매뉴얼을 따른다.}
\end{frame}

\begin{frame}{6. 설치되어 있어도 꺼져 있을 수 있다}
\small 상세 화면에 \textbf{Enable}이 보이면 활성화 상태를 확인한다.
\vscpicture[trim=262.5bp 356.25bp 0bp 52.5bp,clip]{06-enable-workspace}
\vscnote{Enable 오른쪽 화살표 → \textbf{Enable (Workspace)}를 선택한다.\par
이 실습 창에서 활성화한다. 이미 활성화되어 있으면 이 단계는 건너뛴다.}
\end{frame}

\begin{frame}{7. VaporView도 같은 방법으로 확인}
\small \texttt{@installed vaporview} 검색 → 상세 화면 → 활성화 상태 확인.
\vscpicture[trim=262.5bp 356.25bp 0bp 52.5bp,clip]{14-vaporview-enable}
\vscnote{VaporView는 VCD에 저장된 시간을 파형으로 보여주는 확장이다.\par
설치·활성화가 끝나면 \textbf{Explorer} 아이콘으로 돌아간다.}
\end{frame}

\begin{frame}[label=vsc01-truth]{8. 실행 전에 진리표를 먼저 쓴다}
\begin{center}
\renewcommand{\arraystretch}{1.7}
\begin{tabular}{cc|ccc}
$a$ & $b$ & $x$ (AND) & $y$ (OR) & $z$ (XOR)\\\hline
0 & 0 & 0 & 0 & 0\\
1 & 0 & 0 & 1 & 1\\
0 & 1 & 0 & 1 & 1\\
1 & 1 & 1 & 1 & 0
\end{tabular}
\end{center}
\vscnote{AND: 둘 다 1이면 1. OR: 하나 이상 1이면 1.\par
XOR: 두 입력이 서로 다르면 1.\par
이 실습의 입력 순서는 원본과 같은 \textbf{00 → 10 → 01 → 11}이다.}
\end{frame}

\begin{frame}{9. RTL 파일을 연다}
\small \textbf{COMMON → rtl → logic\_gate.v}를 차례로 펼치고 파일을 두 번 누른다.
\vscpicture[trim=262.5bp 318.75bp 315bp 52.5bp,clip]{07-rtl}
\vscnote{입력은 a·b, 출력은 x·y·z다. \texttt{assign}은 조합회로의 연결을 나타낸다.\par
아래 상태 표시줄에 \textbf{Verilog}가 보이는지 확인한다.}
\end{frame}

\begin{frame}{10. 테스트벤치에서 회로를 연결한다}
\small \textbf{COMMON → tb → tb\_logic\_gate.sv}를 두 번 누른다.
\vscpicture[trim=262.5bp 228.75bp 0bp 52.5bp,clip]{08-tb-check}
\vscnote{\texttt{dut}는 검사할 회로다. \texttt{check}는 x·y·z를 예상값과 비교한다.\par
다르면 FAIL로 중단한다. 코드 진단에 빨간 줄이 없다는 것만으로 검사가 끝난 것은 아니다.}
\end{frame}

\begin{frame}{11. 입력 자극·저장·종료를 읽는다}
\small 테스트벤치를 아래로 내려 \texttt{initial} 부분을 본다.
\vscpicture[trim=262.5bp 67.5bp 0bp 247.5bp,clip]{09-tb-stimulus}
\vscnote{\texttt{\#20}은 20ns가 아니라 \textbf{20 × 10ns = 200ns}다.\par
VCD를 저장하고 네 조합 검사를 마친 \textbf{800ns}에 종료한다.}
\end{frame}

\begin{frame}[label=vsc01-run]{12. 수정했다면 저장부터 한다}
\small
\begin{enumerate}
\setlength{\itemsep}{0.22cm}
\item 코드 탭을 누르고 \textbf{File → Save}(\textbf{Ctrl+S})를 선택한다. 여러 파일을 고쳤다면 \textbf{File → Save All}을 사용한다.
\item 코드 탭 이름 옆의 미저장 표시가 사라졌는지 확인한다.
\item \textbf{Terminal → Run Task...}를 선택한다. 창이 좁으면 상단 \textbf{…} 안에 Terminal이 들어 있다.
\end{enumerate}
\vscpicture[height=3.5cm,trim=210bp 315bp 277.5bp 0bp,clip]{10-run-task-menu}
\end{frame}

\begin{frame}{13. 논리 게이트 시뮬레이션 작업 선택}
\small 목록에서 \textbf{LAB1 01: Simulate logic gates}를 누른다.
\vscpicture[height=3.5cm,trim=217.5bp 423.75bp 217.5bp 0bp,clip]{11-task-picker}
\vscnote{파일을 따로 지정하지 않아도 작업이 RTL과 테스트벤치를 함께 사용한다.\par
같은 작업을 \textbf{Terminal → Run Build Task...}(\textbf{Ctrl+Shift+B})로도 실행할 수 있다.\par
작업이 없으면 3단계로 돌아가 지정된 workspace를 열었는지 확인한다.}
\end{frame}

\begin{frame}{14. 완료 로그를 끝까지 확인}
\small 아래 터미널의 작업이 끝날 때까지 기다린다. 결과는 로그 파일에도 저장된다.
\vscpicture[trim=262.5bp 150bp 0bp 217.5bp,clip]{12-simulation-log}
\vscnote{확인할 문구: \texttt{PASS: all 4 logic-gate input combinations}.\par
\textbf{200·400·600·800ns의 PASS 네 줄}과 \textbf{800ns 종료}를 함께 확인한다.}
\end{frame}

\begin{frame}{15. 로그와 파형이 생긴 위치}
\small
\begin{enumerate}
\setlength{\itemsep}{0.18cm}
\item Explorer에서 \textbf{01\_LOGIC\_GATES → build → vscode}를 펼친다.
\item \textbf{simulation.log}를 두 번 누르면 방금 실행한 전체 로그를 읽을 수 있다.
\item \textbf{logic\_gate.vcd}가 파형 파일이다. 다음 단계에서 연다.
\end{enumerate}
\vscpicture[height=3.25cm,trim=33.75bp 345bp 637.5bp 52.5bp,clip]{12-simulation-log}
\vscnote{\texttt{run-...} 폴더에는 컴파일·연결·실행 로그가 따로 남는다.\par
실패했다면 예전 캡처 대신 이번 실행의 오류와 원인을 확인한다.}
\end{frame}

\begin{frame}[label=vsc01-wave]{16. VCD가 글자로 열리면}
\small 파일은 정상이어도 \textbf{Text Editor}로 열면 파형 대신 내용이 보인다.
\vscpicture[trim=258.75bp 225bp 195bp 37.5bp,clip]{13-vcd-as-text}
\vscnote{이 화면에서 값을 직접 고치지 않는다. 다음 단계처럼 \textbf{VaporView}로 다시 연다.}
\end{frame}

\begin{frame}{17. Open With... 선택}
\small Explorer의 \textbf{logic\_gate.vcd}를 마우스 오른쪽 버튼으로 누른다.
\vscpicture[trim=33.75bp 150bp 513.75bp 105bp,clip]{15-open-with}
\vscnote{메뉴에서 \textbf{Open With...}를 누른다.\par
Open in Integrated Terminal과 구분한다.}
\end{frame}

\begin{frame}{18. VaporView로 연다}
\small 편집기 선택 목록에서 \textbf{VaporView}를 누른다.
\vscpicture[trim=217.5bp 457.5bp 217.5bp 0bp,clip]{16-select-vaporview}
\vscnote{VaporView가 목록에 없으면 5--7단계에서 설치·활성화를 다시 확인한다.\par
파일 탭에 파형 아이콘과 VaporView 표시가 나타나는지 확인한다.}
\end{frame}

\begin{frame}{19. 빈 파형 화면에서 신호 목록 열기}
\small 처음에는 신호가 하나도 없을 수 있다. 화면 안의 \textbf{Netlist View}를 누른다.
\vscpicture[trim=258.75bp 390bp 0bp 37.5bp,clip]{17-empty-wave}
\vscnote{왼쪽에 Vaporview: Netlist가 나타난다.\par
\textbf{tb\_logic\_gate} 왼쪽의 삼각형을 눌러 계층을 펼친다.}
\end{frame}

\begin{frame}{20. a·b·x·y·z를 추가한다}
\small \textbf{a, b, x, y, z}를 차례로 두 번 누른다. 오른쪽 파형 목록에 한 줄씩 생긴다.
\vscpicture[height=4.1cm,trim=33.75bp 360bp 637.5bp 22.5bp,clip]{18-netlist}
\vscnote{\texttt{checked}는 검사 횟수이며 회로 출력이 아니다. 이번 파형에는 넣지 않는다.\par
\texttt{dut} 아래에도 같은 회로 신호가 있으므로 중복해서 추가하지 않는다.}
\end{frame}

\begin{frame}{21. 파형이 평평하면 시간 범위를 확인}
\small 파형 위 도구 모음의 \textbf{맨 왼쪽 돋보기(Zoom to Fit)}를 누른다.
\vscpicture[trim=258.75bp 341.25bp 0bp 52.5bp,clip]{20-wave-ps}
\vscnote{확대된 시작 구간만 보다가 전체 구간을 보게 되면 입력 변화가 나타난다.\par
화면의 \texttt{200000 ps}는 \textbf{200ns}와 같다. 다음 단계에서 단위를 바꾼다.}
\end{frame}

\begin{frame}{22. 시간 단위를 ns로 표시}
\small 오른쪽 아래 상태 표시줄의 \textbf{Time Units: ps}를 누르고 \textbf{ns}를 선택한다.
\vscpicture[trim=217.5bp 487.5bp 217.5bp 0bp,clip]{21-select-ns}
\vscnote{파형 데이터 자체를 바꾸는 것이 아니라 표시 단위만 바꾼다.\par
로그의 시간과 파형의 시간을 같은 단위로 비교한다.}
\end{frame}

\begin{frame}{23. 입력 10 구간의 값을 읽는다}
\small 200--400ns 사이 파형을 한 번 누른다. 세로 커서와 신호별 값이 나타난다.
\vscpicture[trim=258.75bp 360bp 0bp 52.5bp,clip]{23-wave-10}
\vscnote{실제 캡처의 커서는 약 296.7ns다. \textbf{a=1, b=0, x=0, y=1, z=1}.\par
전환 경계 대신 구간 안쪽을 읽는다. 다른 세 구간도 같은 방법으로 확인한다.}
\end{frame}

\begin{frame}{24. 입력 11에서는 XOR가 0이다}
\small 600--800ns 사이를 눌러 마지막 입력 조합을 확인한다.
\vscpicture[trim=258.75bp 360bp 0bp 52.5bp,clip]{22-wave-11}
\vscnote{약 699.8ns에서 \textbf{a=1, b=1, x=1, y=1, z=0}.\par
OR와 XOR의 차이가 드러나는 구간이다. 이 차이를 레포트에서 설명한다.}
\end{frame}

\begin{frame}{25. 파형과 진리표를 한 줄씩 비교}
\small\renewcommand{\arraystretch}{1.6}
\begin{center}
\begin{tabular}{c|cc|ccc}
시간 구간(ns) & a & b & x & y & z\\\hline
0--200 & 0 & 0 & 0 & 0 & 0\\
200--400 & 1 & 0 & 0 & 1 & 1\\
400--600 & 0 & 1 & 0 & 1 & 1\\
600--800 & 1 & 1 & 1 & 1 & 0
\end{tabular}
\end{center}
\vscnote{로그는 각 구간의 끝에서 값을 검사하고, 파형은 구간 전체의 변화를 보여준다.\par
원본과 같은 입력 순서이며 \textbf{800ns 이후로 테스트가 더 진행된 것은 아니다.}\par
화면 맞춤 때문에 종료 시각 바깥까지 축이 보일 수 있다.}
\end{frame}

\begin{frame}[label=vsc01-report]{26. 실험 전 레포트에 기록할 내용}
\small
\begin{enumerate}
\setlength{\itemsep}{0.18cm}
\item 회로 목적, a·b 입력과 x·y·z 출력, 실행 전에 작성한 진리표.
\item RTL의 AND·OR·XOR 연결과 테스트벤치의 네 입력 조합.
\item 사용한 workspace·시뮬레이터 버전·작업 이름, 시간 단위와 200ns 간격·800ns 종료 이유.
\item PASS 네 줄과 최종 완료 문구를 포함한 \textbf{자신의 실행 로그 캡처}.
\item a·b·x·y·z, 시간축, 커서 값이 보이는 \textbf{자신의 파형 캡처}.
\item 네 구간별 예상값·실제값 비교, OR와 XOR가 달라지는 이유, 오류가 있었다면 수정·재실행 과정.
\end{enumerate}
\vscnote{이 자료의 캡처를 자신의 실험 결과로 제출하지 않는다.\par
코드를 저장소에서 가져와도 구조와 결과를 설명할 수 있어야 한다.}
\end{frame}

\begin{frame}{27. 수정 후에는 다시 실행한다}
\small
\begin{enumerate}
\setlength{\itemsep}{0.23cm}
\item RTL 또는 테스트벤치를 수정했다면 \textbf{Save / Save All}.
\item \textbf{Run Task... → LAB1 01: Simulate logic gates} 재실행.
\item 이번 실행의 PASS·종료 시각과 로그를 다시 확인.
\item VaporView 탭 오른쪽 위의 \textbf{새로고침(원형 화살표)}을 누르거나 VCD를 다시 연다.
\item 동일한 신호·시간 구간을 비교하고 새 캡처를 레포트에 넣는다.
\end{enumerate}
\vscnote{소스만 고치고 이전 파형을 보면 수정 결과를 확인할 수 없다.\par
실행 실패 후에는 이전 파형을 성공 결과로 사용하지 않는다.}
\end{frame}

\begin{frame}{28. 막혔을 때 돌아갈 곳}
\small\renewcommand{\arraystretch}{1.35}
\begin{tabular}{@{}p{3.5cm}p{6.8cm}@{}}
\textbf{보이는 문제} & \textbf{확인할 것}\\\hline
시뮬레이션 작업 없음 & 지정한 .code-workspace를 열었는지(3단계).\\
도구를 찾을 수 없음 & Vivado bin의 PATH 또는 VIVADO\_BIN 설정 후 VS Code 다시 열기.\\
코드가 Plain Text & slang 확장 활성화·파일 확장자(5--6단계).\\
모듈을 찾지 못하는 진단 & sources.f와 .slang 설정 확인, 명령 팔레트의 slang: Restart Language Server.\\
VCD가 글자로 보임 & Open With... → VaporView(16--18단계).\\
파형 없음·평평한 파형 & PASS 확인 → 신호 추가 → Zoom to Fit.\\
FAIL 또는 컴파일 오류 & 이번 실행의 로그·소스 수정 → 저장 → 재실행.
\end{tabular}
\end{frame}

\begin{frame}{29. 이제 Vivado GUI로 넘어간다}
\small
\begin{itemize}
\setlength{\itemsep}{0.23cm}
\item 네 입력 조합의 PASS, 800ns 종료, 파형·진리표의 일치를 확인했다.
\item 실험 전 레포트에 실행 방법·코드 설명·실제 로그·파형 해석을 기록했다.
\item 다음은 레거시 원본의 Vivado 절차다. 프로젝트를 만들고 같은 회로·입력을 Vivado에서도 비교한다.
\item 실험 후에는 실제 FPGA에 비트스트림을 기록하고 동작 사진·영상·GitHub 링크를 실험 후 레포트에 정리한다.
\end{itemize}
\vspace{0.35cm}
\navlink{legacy-01-vivado}{▶ 레거시 Vivado 실습 시작}\quad
\navlink{legacy-01-vivado-contents}{Vivado 세부 목차}
\vscnote{공식 참고: \href{https://code.visualstudio.com/docs/editing/workspaces/workspaces}{VS Code workspace} ·
\href{https://code.visualstudio.com/docs/debugtest/tasks}{Tasks} ·
\href{https://github.com/Lramseyer/vaporview}{VaporView}}
\end{frame}
\endgroup
